\scanpage{35}
l'on doit avoir, si $\unsure{x|a}\in V\cap T'(\unsure{V|W})$, $R_K\subset G_{x\,K}$. L'inverse est vrai, car comme $T_K$ est fid. plat \unsure{par|sur} $K$, on voit par descente que les \unsure{pts} (\unsure{à coeff. ds des|à coeff. dans des} $K$-alg. quelconques $\Lambda$) $s.x$ $(s\in T_K(\Lambda))$ proviennent d'un pt bien déterminé de $V_K$. Or soit \deleted{\unsure{$\xi$}}
\[
S_{\unsure{K|k}}=\overline{\unsure{\varphi|p}(R_K)}\qquad p\colon G_K\unsure{\mathrel{-}|\longrightarrow}G,
\]
et soit $H$ le ss-gpe algébrique \inserted{fermé} de $G$ engendré par $S$. On a donc
\[
x\in V\cap T'(V)\quad\Longleftrightarrow\quad x\in V_{\deleted{\unsure{x|0}}}^{H}\quad \deleted{\unsure{H|\forall}\unsure{\subset|\in}\unsure{G}\ \unsure{H_x|H}\ \text{\unsure{tous|tels}}}.
\]

Par suite, pour que \deleted{les invariants} pour tout $V$, on ait $V\cap T'(V)=V^G$, i.e. que le foncteur $\mathrm{Rep}(G)=\unsure{\mathcal M|M}\longrightarrow$ (structures vectorielles mixtes pour $L/K/k)$ soit pl. fid, il \unsure{n. f. et s.|f. et s.} que pour $\unsure{\forall|H}\,V$, \deleted{les invariants} on ait
\[
V^H=V^G.
\]

Supposons $P$ quelconque. On a l'\unsure{hom|hom.} canonique de bitorseurs
\[
\deleted{\unsure{\underline{\mathrm{Isom}}|{\mathrm{Isom}}}_{\unsure{\otimes}}(\unsure{T'|T},T_K))}
\]
\diagram{\begin{tikzcd}[ampersand replacement=\&,row sep=small]
P\mathbin{\unsure{\times|\otimes}}_K P\arrow[r] \& G_K\\
(v,u)\arrow[r,mapsto] \& v^{-1}u
\end{tikzcd}}
\deleted{\unsure{induisant}}
\[
\deleted{P_L^{\unsure{a|G}}\times P_L\mathrel{\unsure{-|\longrightarrow}}G_L}
\]
donc $\varphi\colon\mathrm{Spec}\,L\longrightarrow P$ (un $\mathrm{Spec}\,K$-\unsure{mor.|m.}, bien sûr) s'en déduit \deleted{de, grâce à l'isom $P_L\simeq G_L$}
\[
\mathrm{Spec}(L)\times_{\mathrm{Spec}\,K}\mathrm{Spec}(L)\mathrel{\unsure{-|\longrightarrow}}
P\times_{\mathrm{Spec}\,\unsure{K|k}}P\longrightarrow G_K\,;
\]
l'image fermée $\deleted{\unsure{P|R}}R_K$ est donc définie intrinsèquement. On vérifie encore (par descente \unsure{platisation|fidèlement|plate}) que \deleted{\unsure{p}} l'on a, pour tt $G$-module $V$
\[
\deleted{\unsure{x\in}}V\cap T'(V)=\{\,x\in V,\mid R_K\subset G_{x\,K}\,\}.
\]
Définissant $H$ comme ci-dessus, on trouve donc
\[
V\cap T'(V)=V^H,
\]
et la \unsure{$\overline{\mathrm m}$|même} conclusion que plus haut ....
% paper watermark, sideways
{\centering \textsc{Extra Strong}\\
\textit{Ma\unsure{in|nn}\unsure{ol|on}\unsure{uy|ny|y} 1808}\par}
% archive sheet number: 34
